%%%%%%%%%%%%%%%%%%%%%%%%%%%%%%%%%%%%%%%%%
% CPS Online Quiz System Poster
% Based on standard academic poster templates (Beamerposter)
% Optimized for Overleaf
%%%%%%%%%%%%%%%%%%%%%%%%%%%%%%%%%%%%%%%%%

\documentclass[final]{beamer}

% ==================== PACKAGES ====================
\usepackage[size=a1,orientation=landscape,scale=1.4]{beamerposter} 
%\usetheme{gemini} % If gemini theme is not available, standard beamer works, but we will style manually below to ensure it runs without external theme files.
\usecolortheme{beaver} % Red/Grey theme (Clean academic look)
\usepackage[utf8]{inputenc}
\usepackage{graphicx}
\usepackage{booktabs}
\usepackage{tikz}
\usepackage{pgfplots}
\usepackage{anyfontsize}
\usepackage{tcolorbox}
\usepackage{palatino} % Professional font

% ==================== MANUAL STYLING (To ensure it looks like the reference) ====================
% Define colors similar to the CCGRID poster (Blue/Grey/White)
\definecolor{VUblue}{RGB}{26, 59, 92} % Dark Blue header
\definecolor{VUgold}{RGB}{255, 204, 0} % Accent
\definecolor{BlockBG}{RGB}{255, 255, 255}
\definecolor{BlockTitle}{RGB}{26, 59, 92}

\setbeamercolor{headline}{bg=VUblue,fg=white}
\setbeamercolor{footline}{bg=VUblue,fg=white}
\setbeamercolor{title in headline}{fg=white}
\setbeamercolor{author in headline}{fg=white}
\setbeamercolor{institute in headline}{fg=lightgray}

% Style the blocks to look like the uploaded reference
\setbeamercolor{block title}{bg=VUblue,fg=white}
\setbeamercolor{block body}{bg=white,fg=black}
\setbeamertemplate{block begin}{
  \begin{tcolorbox}[colback=white,colframe=VUblue,arc=4pt,title=\insertblocktitle, boxrule=1mm, titlerule=0mm, toptitle=2mm, bottomtitle=2mm]
}
\setbeamertemplate{block end}{\end{tcolorbox}\vspace{1cm}}

% ==================== METADATA ====================
\title{\huge Online Quiz System: A Java OOP Implementation}
\author{\large Hanwen Zhang (1308159) \& Zimo Qin}
\institute{\large Wenzhou-Kean University | CPS 2231 W02}
\date{\today}

% ==================== DOCUMENT ====================
\begin{document}
\begin{frame}[t]

% --- HEADER ---
\begin{beamercolorbox}[wd=\paperwidth]{headline}
  \begin{columns}[T]
    \begin{column}{.02\paperwidth}\end{column}
    \begin{column}{.70\paperwidth}
      \vskip2cm
      \centering
      \usebeamerfont{title in headline}{\textbf{\huge Online Quiz System: A Java OOP Implementation}}\\[1ex]
      \usebeamerfont{author in headline}{\large Hanwen Zhang (1308159) \& Zimo Qin | CPS 2231 W02}\\[1ex]
      \usebeamerfont{institute in headline}{Wenzhou-Kean University}\\[1cm]
    \end{column}
    \begin{column}{.25\paperwidth}
      \vskip1.5cm
      % PLACEHOLDER FOR UNIVERSITY LOGO
      % \includegraphics[width=0.7\linewidth]{logo.png} 
      \begin{tikzpicture}
        \node[draw=white, text=white, very thick, inner sep=10pt] {\LARGE LOGO HERE};
      \end{tikzpicture}
      \vskip1cm
    \end{column}
    \begin{column}{.02\paperwidth}\end{column}
  \end{columns}
\end{beamercolorbox}

% --- COLUMNS ---
\begin{columns}[T]

  % ==================== COLUMN 1 ====================
  \begin{column}{.32\paperwidth}
    
    \begin{block}{1. System Overview}
      This project presents a robust **Online Quiz System** built using Java Object-Oriented Programming (OOP) principles. The system encapsulates question display, answer verification, and input validation into a modular architecture.
      
      \begin{itemize}
          \item \textbf{Core Philosophy:} Designed using abstract classes, inheritance, and method overriding to ensure extensibility.
          \item \textbf{Key Functionality:} Manages the full quiz process, including input detection, score calculation, and result display.
      \end{itemize}
    \end{block}

    \begin{block}{2. System Architecture}
      The application follows a hierarchical class structure centered around the abstract \texttt{Question} class.
      
      \vspace{0.5cm}
      
      % --- DIAGRAM PLACEHOLDER (Use actual UML image here) ---
      \begin{center}
      \begin{tikzpicture}
        % Simple UML Representation drawn in LaTeX
        \node (Abstract) [draw, rectangle, fill=blue!10, text width=6cm, align=center] {\textbf{<<Abstract>>}\\ Question\\ \textit{+ checkAnswer()}};
        \node (MC) [draw, rectangle, below left of=Abstract, node distance=5cm, xshift=-2cm, text width=4cm, align=center] {MultipleChoice\\ Question};
        \node (TF) [draw, rectangle, below right of=Abstract, node distance=5cm, xshift=2cm, text width=4cm, align=center] {TrueFalse\\ Question};
        \node (Manager) [draw, rectangle, fill=green!10, above of=Abstract, node distance=4cm, text width=6cm, align=center] {\textbf{QuizManager}\\ \textit{Input Validation \& Scheduling}};
        
        \draw[->, thick] (MC) -- (Abstract);
        \draw[->, thick] (TF) -- (Abstract);
        \draw[dashed, thick] (Manager) -- (Abstract);
      \end{tikzpicture}
      \end{center}
      \vspace{0.5cm}

      \textbf{Components:}
      \begin{itemize}
          \item \textbf{Inheritance Tree:} \texttt{MultipleChoice} and \texttt{TrueFalse} inherit from \texttt{Question}.
          \item \textbf{Management:} \texttt{QuizManager} schedules the interaction between \texttt{Student} entities and question logic.
      \end{itemize}
    \end{block}

  \end{column}

  % ==================== COLUMN 2 ====================
  \begin{column}{.32\paperwidth}

    \begin{block}{3. Implementation Details}
      \textbf{The Abstract Question Class} serves as the blueprint for all question types.
      
      \vspace{0.5em}
      \textbf{Polymorphism in Action:}
      \begin{itemize}
          \item \textbf{Multiple Choice:} Overrides \texttt{display()} to automatically generate A/B/C/D options using character offsets.
          \item \textbf{True/False:} Implements strict logic where "A" maps to True and "B" maps to False.
      \end{itemize}

      \vspace{1em}
      \begin{tcolorbox}[colback=gray!10, title=Code Snippet: Polymorphism]
\begin{verbatim}
// Abstract Definition
public abstract boolean checkAnswer(String input);

// Implementation in TrueFalseQuestion
@Override
public boolean checkAnswer(String answer) {
    // Validates A/B input logic
    return answer.equalsIgnoreCase(correctAnswer);
}
\end{verbatim}
      \end{tcolorbox}
    \end{block}

    \begin{block}{4. Workflow Diagram}
      The system ensures valid data entry through a cyclical check within \texttt{QuizManager}.
      
      \vspace{1cm}
      % --- FLOWCHART PLACEHOLDER ---
      \begin{center}
      \begin{tikzpicture}[node distance=2.5cm]
        \node (start) [draw, circle, fill=green!20] {Start};
        \node (input) [draw, rectangle, below of=start] {User Input};
        \node (validate) [draw, diamond, aspect=2, below of=input] {Valid Format?};
        \node (error) [draw, rectangle, right of=validate, xshift=2cm, fill=red!10] {Throw Exception};
        \node (process) [draw, rectangle, below of=validate] {Process Answer};
        
        \draw[->] (start) -- (input);
        \draw[->] (input) -- (validate);
        \draw[->] (validate) -- node[right] {No} (error);
        \draw[->] (error) |- (input);
        \draw[->] (validate) -- node[left] {Yes} (process);
      \end{tikzpicture}
      \end{center}
      \vspace{0.5cm}
      
      \textit{Figure 2: The logic flow for handling student input and exceptions.}
    \end{block}

  \end{column}

  % ==================== COLUMN 3 ====================
  \begin{column}{.32\paperwidth}

    \begin{block}{5. Robustness \& Validation}
       A critical feature is the \texttt{validateInput} mechanism within the \texttt{QuizManager}.
       
       \begin{itemize}
           \item \textbf{Exception Handling:} Uses \texttt{InvalidInputException} to handle erroneous user data.
           \item \textbf{Empty Input:} Rejects empty strings immediately to prevent null pointers.
           \item \textbf{Type Checking:} For True/False questions, inputs other than "A" or "B" (case-insensitive) trigger a retry loop.
           \item \textbf{Thread Safety:} The \texttt{Student} class uses the \texttt{synchronized} keyword for \texttt{addScore}, enabling future expansion to multi-threaded environments.
       \end{itemize}
    \end{block}

    \begin{block}{6. Conclusion}
      We successfully implemented an OOP-based Quiz System that prioritizes modularity and error resilience.
      
      \vspace{0.5em}
      \textbf{Key Achievements:}
      \begin{itemize}
          \item \textbf{Reliability:} Secured via custom exception handling.
          \item \textbf{Scalability:} Abstract design allows easy addition of new question types (e.g., Fill-in-Blank) without breaking existing code.
          \item \textbf{Maintainability:} Clear separation of logic (\texttt{Question}), data (\texttt{Student}), and control (\texttt{QuizManager}).
      \end{itemize}
    \end{block}

    \begin{block}{References}
      \small
      \begin{enumerate}
          \item Horstmann, C. S. (2018). \textit{Core Java Volume I--Fundamentals}. Prentice Hall.
          \item Oracle Documentation. (2023). \textit{Java Platform, Standard Edition & Java Development Kit Version 17 API Specification}.
      \end{enumerate}
    \end{block}

  \end{column}

\end{columns}
\end{frame}
\end{document}